% GENERATED by scripts/build-lecture-slides.py; do not edit. Portable profile: XeTeX, theme and fonts from the core.
\makeatletter\IfFileExists{core/beamer/beamerthememx.sty}{\def\mxcorepath{core/}\def\input@path{{core/beamer/}{core/}}}{\def\mxcorepath{../core/}\def\input@path{{../core/beamer/}{../core/}}}\makeatother
\documentclass[aspectratio=169,10pt,t]{beamer}
\usetheme{mx}
\usepackage{graphicx,booktabs,mathtools,array,adjustbox}
\usepackage[all]{xy}
\providecommand{\actualtextwidth}{\textwidth}
\providecommand{\veccor}[2]{\left(\begin{matrix}#1\\#2\end{matrix}\right)}
\DeclarePairedDelimiter{\VecNorm}{\lVert}{\rVert}
\providecommand{\ETC}{\ldots}
\usepackage[french]{babel}
\setlength{\parskip}{.45em}
\hypersetup{pdfauthor={},pdfsubject={Lecture slides},pdfkeywords={},pdftitle={Primitives et intégrales}}
\providecommand{\N}{\mathbb{N}}
\providecommand{\Z}{\mathbb{Z}}
\providecommand{\D}{\mathbb{D}}
\providecommand{\Q}{\mathbb{Q}}
\providecommand{\R}{\mathbb{R}}
\providecommand{\C}{\mathbb{C}}
\providecommand{\vect}[1]{\overrightarrow{#1}}
\providecommand{\vi}{\vec{\imath}}
\providecommand{\vj}{\vec{\jmath}}
\providecommand{\vk}{\vec{k}}
\providecommand{\norm}[1]{\left\lVert #1\right\rVert}
\providecommand{\vc}[2]{\begin{pmatrix}#1\\#2\end{pmatrix}}
\providecommand{\deter}[2]{\det\left(#1,#2\right)}
\providecommand{\abs}[1]{\left\lvert #1\right\rvert}
\providecommand{\dd}{\,\mathrm{d}}
\providecommand{\e}{\mathrm{e}}
\providecommand{\Card}{\operatorname{Card}}
\providecommand{\E}{\operatorname{E}}
\providecommand{\V}{\operatorname{V}}
\providecommand{\degre}{^{\circ}}
\providecommand{\un}[1]{\,\mathrm{#1}}
\providecommand{\sci}[2]{#1\times10^{#2}}
\providecommand{\pgcd}{\operatorname{PGCD}}
\providecommand{\ppcm}{\operatorname{PPCM}}
\title{Primitives et intégrales}
\subtitle{Cours et exemples}
\date{}
\author{}
\begin{document}
\mxtitleframe
\begin{frame}[allowframebreaks]{Primitives}
Soit \(f\) une fonction définie sur un intervalle \(I\) non-vide.
            Une fonction \(F\) est une \emph{primitive} de \(f\) sur \(I\) si \(F\) est dérivable sur \(I\) et si, 
            pour tout \(x\in I\), \(F'(x)=f(x)\).
\end{frame}
\begin{frame}[allowframebreaks]{Primitives (2)}
Si une fonction \(f\) admet une primitive \(F\), alors elle en admet une infinité.
            En fait, toutes les fonctions de la forme \(F+c\) où \(c\in\mathbb{R}\) sont aussi des primitives de \(f\).

            L'ensemble des primitives de \(f\) est noté par \(\displaystyle{\int f(x) ~\mathrm{d}x}\).
            example
\end{frame}
\begin{frame}[allowframebreaks]{Primitives (3)}
\begin{itemize}

                        \item \(f:x\mapsto 2x\) admet \(F:x\mapsto x^2\) comme primitive sur \(\mathbb{R}\).
                        
\item \(g:x\mapsto \dfrac{1}{2\sqrt{x}}\) admet \(G:x\mapsto \sqrt{x} + 2\) 
                              comme primitive sur \(]0,+\infty[\).

\end{itemize}

            example
\end{frame}
\begin{frame}[allowframebreaks]{Primitives}
Soit \(f\) une fonction définie sur un intervalle \(I\) et \(F\) une primitive de \(f\) sur \(I\).
\end{frame}
\begin{frame}[allowframebreaks]{Primitives (2)}
\begin{enumerate}
\setcounter{enumi}{0}\item La fonction \(\lambda f\) admet la fonction \(\lambda F\) comme primitive sur \(I\).
            
\item Soit \(g\) une fonction définie sur un intervalle \(I\) et \(G\) une primitive de \(g\) sur \(I\).
                  Alors la fonction \(f+g\) admet la fonction \(F+G\) comme primitive sur \(I\).      
            
\item Soit \(F_1\) une autre primitive de \(f\) sur \(I\).
                  Alors la fonction \(F - F_1\) est constante.

\end{enumerate}
\end{frame}
\begin{frame}[allowframebreaks]{Primitives (3)}
\begin{enumerate}
\setcounter{enumi}{3}\item Soient \(x_0\in I\) et \(y_0\in \mathbb{R}\).
                  Alors \(f\) admet une unique primitive \(F_0\) sur \(I\) vérifiant \(F_0 (x_0) = y_0\).

\end{enumerate}
\end{frame}
\begin{frame}[allowframebreaks]{Primitives (4)}
example
                  Cherchons la primitive de la fonction \(f:x\mapsto 3 x^2\) qui s'annule en \(1\).
                  La fonction \(F:x\mapsto x^3\) est dérivable et sa dérivée est \(F':x\mapsto 3x^2\).
                  Comme \(F(1)=1\), alors on peut prendre \(F_0 : x\mapsto x^3 - 1\).
                  Ainsi, \(F_0' (x) = 3x^2\) et \(F_0 (1) = 1 - 1 = 0\).
            example
\end{frame}
\begin{frame}[allowframebreaks]{Primitives}
minipage0.45
\end{frame}
\begin{frame}[allowframebreaks]{Primitives (2)}
\footnotesize
\begin{center}\begin{tabular}{ll}
Fonctions & Primitives \\ {}
\(0\) & \(k\) \\ {}
\(1\) & \(x + k\) \\ {}
\(\dfrac{1}{x^2}\) & \(-\dfrac{1}{x}+k\) \\ {}
\(\dfrac{1}{\sqrt{x}}\) & \(2\sqrt{x}+k\) \\ {}
\(x^n \quad (n\neq -1)\) & \(\dfrac{x^{n+1}}{n+1} + k\) \\ {}
\(\cos x\) & \(\sin x + k\) \\ {}
\(\sin x\) & \(-\cos x +k\) \\ {}
\(\mathrm{e}^{x}\) & \(\mathrm{e}^{x} + k\) \\ {}
\(\dfrac{1}{x}\) & \(\ln |x| + k\) \\ {}
\end{tabular}\end{center}
\end{frame}
\begin{frame}[allowframebreaks]{Primitives (3)}
minipage
      
            minipage0.45
                  Notons que la recherche de primitives est en général difficile.
                  Bien que toute fonction continue admette une primitive, 
                  il existe des fonctions dont on ne 
                  sait pas donner des primitives sous forme de fonctions usuelles.

En voici quelques exemples :
\end{frame}
\begin{frame}[allowframebreaks]{Primitives (4)}
\begin{itemize}

							\item \(x \mapsto \mathrm{e}^{-x^2}\)
							
\item \(x \mapsto \sqrt{x^4 + 1}\)
							
\item \(x \mapsto (\ln x)^{-1}\)
							
\item \(x \mapsto \sqrt{x} \cdot \sin (x) \)
							
\item \ldots{}

\end{itemize}

            minipage
\end{frame}
\begin{frame}[allowframebreaks]{Primitives}
Soient \(g\) une fonction dérivable sur un intervalle \(I\) et \(f\) une fonction dérivable sur un intervalle
            contenant \(g(I)\).
            Comme la recherche d'une primitive est l'opération inverse de la dérivation, on a :
                  \[\begin{array}{rcl}
                        \int g'(x) \cdot f'\left( g(x) \right) ~\mathrm{d}x & = & f\left( g(x) \right) + C       \\
                        \int f'(x) \cdot f^n (x) ~\mathrm{d}x               & = & \dfrac{1}{n+1} f^{n+1}(x) + C
                        \quad (n \in \mathbb{N})
                  \end{array}\]
\end{frame}
\begin{frame}[allowframebreaks]{Primitives (2)}
En particulier, pour une fonction \(u\) dérivable sur \(I\) (et, selon le cas, \(u(x)\neq 0\) ou \(u(x)>0\) sur \(I\)) :
                  \[\begin{array}{rcl}
                        \int \dfrac{u'(x)}{u(x)} ~\mathrm{d}x        & = & \ln |u(x)| + C ,      \\
                        \int u'(x) \, \mathrm{e}^{u(x)} ~\mathrm{d}x      & = & \mathrm{e}^{u(x)} + C ,     \\
                        \int \dfrac{u'(x)}{\sqrt{u(x)}} ~\mathrm{d}x & = & 2 \sqrt{u(x)} + C ,   \\
                        \int \dfrac{u'(x)}{u^2(x)} ~\mathrm{d}x      & = & -\dfrac{1}{u(x)} + C .
                  \end{array}\]

example
\end{frame}
\begin{frame}[allowframebreaks]{Primitives (3)}
\footnotesize
\begin{enumerate}
\setcounter{enumi}{0}\item On veut trouver une primitive de \(a(x) = 2x \cdot \cos (x^2)\).
                              On constate que \(2x\) est en fait la dérivée de \(x^2\).
                              Cette fonction est donc de la forme \(u'(x) \cdot \cos u(x)\).
                              Il nous reste donc à trouver la primitive de la fonction \(\cos\),
                              qui n'est autre que la fonction \(\sin\).
                              Ainsi, 
                                    \[\begin{aligned}
                                    a(x) & = (x^2)' \cdot \sin' (x^2)            \\
                                         & = g'(x) \cdot f'\left(g(x)\right)     \,.
                                    \end{aligned}\]
                              Ainsi, la fonction \(A\) définie par \(A(x) = \sin (x^2)\) est une primitive de \(a\).

\end{enumerate}
\end{frame}
\begin{frame}[allowframebreaks]{Primitives (4)}
\begin{enumerate}
\setcounter{enumi}{1}\item On veut trouver une primitive de \(b(x) = \cos x \cdot \sin^4 x\).
                              On commence par remarquer que \(\cos\) est en fait la dérivée de \(\sin\).
                              Ainsi, \(b(x) = \sin' x \cdot \sin^4 x\).
                              La fonction \(B\) définie par \(B(x) = \dfrac{\sin^5 x}{5}\) est une primitive de \(b\).

\end{enumerate}
\end{frame}
\begin{frame}[allowframebreaks]{Primitives (5)}
example
\end{frame}
\begin{frame}[allowframebreaks]{Intégrales}
Soit \(f\) une fonction définie sur un intervalle \(I\) et y admettant une primitive \(F\).
	      Soient \(a\) et \(b\) deux éléments de \(I\).
	      On appelle intégrale de \(f\) entre \(a\) et \(b\) le réel 
	            \[
	            \int_{a}^{b} f(x) ~\mathrm{d}x = \left[ F(x) \right]_{a}^{b} = F(b) - F(a)      \,.
	            \]
\end{frame}
\begin{frame}[allowframebreaks]{Intégrales (2)}
\footnotesize
La variable \(x\) de la définition est muette, 
	              c'est-à-dire qu'on peut la remplacer par n'importe quelle autre variable.
	              Ainsi, \(\int_{a}^{b} f(x) ~\mathrm{d}x = \int_{a}^{b} f(t) \mathrm{d}t\). 
	              La valeur de \(\int_{a}^{b} f(x) ~\mathrm{d}x\) ne dépend pas de la primitive \(F\) choisie.
	              En effet, si \(F_1\) et \(F_2\) sont deux primitives de \(f\), alors 
	              il existe \(C\in\mathbb{R}\) tel que \(F_1 = F_2 + C\).
	              Ainsi, 
	                    \[\begin{aligned}
	                    \left[ F_1 (x) \right]_{a}^{b} & = F_1 (b) - F_1 (a)                    
	                                                     = (F_2 (b) + C ) - (F_2 (a) + C)       \\
	                                                   & = F_2 (b) - F_2 (a)                    
	                                                     = \left[ F_2 (x) \right]_{a}^{b}       \,.
	                    \end{aligned}\]
\end{frame}
\begin{frame}[allowframebreaks]{Intégrales}
Soient \(f\) et \(g\) deux fonctions définies sur un intervalle \(I\) et y admettant des primitives.
	      Soient \(F\) une primitive de \(f\) et \(G\) une primitive de \(g\).
	      Soient \(a\), \(b\) et \(c\) trois éléments de \(I\) et \(\lambda \in \mathbb{R}\).
\end{frame}
\begin{frame}[allowframebreaks]{Intégrales (2)}
\begin{enumerate}
\setcounter{enumi}{0}\item \(\int_{a}^{b} f(x) ~\mathrm{d}x = -\int_{b}^{a} f(x) ~\mathrm{d}x\)
	                  
\item \(\int_{a}^{b} (f+g) (x) ~\mathrm{d}x = \int_{a}^{b} f(x) ~\mathrm{d}x + \int_{a}^{b} g(x) ~\mathrm{d}x\)
	                  
\item \(\int_{a}^{b} \lambda f(x) ~\mathrm{d}x = \lambda \int_{a}^{b} f(x) ~\mathrm{d}x\)

\end{enumerate}
\end{frame}
\begin{frame}[allowframebreaks]{Intégrales (3)}
\begin{enumerate}
\setcounter{enumi}{3}\item \(\int_{a}^{b} f(x) ~\mathrm{d}x = \int_{a}^{c} f(x) ~\mathrm{d}x + \int_{c}^{b} f(x) ~\mathrm{d}x\)
	                  
\item Si \(a \leq b\) et \(f(x) \leq g(x)\) pour tout \(x\in [a,b]\) alors 
	                        \(\int_{a}^{b} f(x) ~\mathrm{d}x \leq \int_{a}^{b} g(x) ~\mathrm{d}x\).            

\end{enumerate}
\end{frame}
\begin{frame}[allowframebreaks]{Intégrales}
Soit \(f\) une fonction définie et continue sur un intervalle \([a,b]\) (\(a<b\)).
\end{frame}
\begin{frame}[allowframebreaks]{Intégrales (2)}
\begin{enumerate}
\setcounter{enumi}{0}\item Si \( f(x) \leq 0\) pour tout \(x\in [a,b]\), alors \(\int_{a}^{b} f(x) ~\mathrm{d}x \leq 0\).
	                  
\item Si \( f(x) \geq 0\) pour tout \( x\in [a,b]\), alors \(\int_{a}^{b} f(x) ~\mathrm{d}x \geq 0\).
	                  
\item \( \left\lvert \int_{a}^{b} f(x) ~\mathrm{d}x \right\rvert \leq \int_{a}^{b} |f(x)| ~\mathrm{d}x \).

\end{enumerate}
\end{frame}
\begin{frame}[allowframebreaks]{Intégrales (3)}
\begin{enumerate}
\setcounter{enumi}{3}\item Si \(m\leq f(x) \leq M\) pour tout \(x\in [a,b]\), 
	                        alors \( m \leq \dfrac{1}{b-a} \int_{a}^{b} f(x) ~\mathrm{d}x \leq M\).

\end{enumerate}
\end{frame}
\begin{frame}[allowframebreaks]{Intégrales (4)}
Le réel \(\mu = \dfrac{1}{b-a} \int_{a}^{b} f(x) ~\mathrm{d}x\) est appelé la \emph{valeur moyenne} de \(f\) sur \([a,b]\) ;
	            le point 4 est l'\emph{inégalité de la moyenne}.
\end{frame}
\begin{frame}[allowframebreaks]{Intégrales}
Soient \(u\) et \(v\) deux fonctions dérivables sur un intervalle non-vide \(I\).
	            Si \(u'v\) admet une primitive sur \(I\), alors on a 
	                  \[
	                        \int u(x) \cdot v'(x) ~\mathrm{d}x = u(x) \cdot v(x) - \int u'(x) \cdot v(x) ~\mathrm{d}x      
	                  \]

ou, avec d'autres notations :
	                  \[
	                        \int u \ \mathrm{d}v = u \cdot v - \int v \ \mathrm{d}u      \,.
	                  \]
\end{frame}
\begin{frame}[allowframebreaks]{Intégrales (2)}
example
	                  On considère la fonction \(f\) définie par \(f(x) = x\cdot \mathrm{e}^{x}\).
	                  On veut trouver la primitive \(F\) de \(f\) qui vérifie \(F(1) = 1\).
	                  
	                  Intégrons par partie en posant
	                        \[\begin{aligned}
	                              u           & = x       & \text{ et } & &      \mathrm{d}v & = \mathrm{e}^{x} ~\mathrm{d}x      \\
	                              \mathrm{d}u & = ~\mathrm{d}x     & \text{    } & &      v           & = \mathrm{e}^{x}          \;.
	                        \end{aligned}\]
\end{frame}
\begin{frame}[allowframebreaks]{Intégrales (3)}
On a alors :
	                        \[\begin{aligned}
	                        \int x \cdot \mathrm{e}^{x} ~\mathrm{d}x & = [x \cdot \mathrm{e}^{x} ] - \int \mathrm{e}^{x} ~\mathrm{d}x       \\
	                                                  & = x \cdot \mathrm{e}^{x} - \mathrm{e}^{x} + C               \;.
	                        \end{aligned}\]
\end{frame}
\begin{frame}[allowframebreaks]{Intégrales (4)}
Posons alors \(F(x) = \mathrm{e}^{x} (x-1) + C\).
	                  On cherche alors la valeur de \(C\) pour que \(F(1) = 1\).
	                  On a \(
	                              F(1) = \mathrm{e}^{1} (1-1) + C = C
	                       \),
	                       ce qui nous donne \(C = 1\).
	                  D'où \(F(x) = \mathrm{e}^{x} (x-1) + 1\).
	            example
\end{frame}
\begin{frame}[allowframebreaks]{Intégrales}
\footnotesize
Soit \(f\) une fonction définie sur un intervalle \(I\) et y admettant une primitive \(F\) 
	      et soit \(\varphi\) une fonction définie et dérivable sur un intervalle \(J\) telle que \(\varphi (J) \subseteq I\).
	      La fonction \(t \mapsto F\left( \varphi (t) \right)\) est dérivable sur \(J\) et a pour dérivée 
	      \(t \mapsto \varphi'(t) \cdot F' \left( \varphi (t) \right)  = \varphi'(t) \cdot f\left(\varphi (t) \right) \).
	      
	      Si de plus \(\varphi\) est une bijection de \(J\) vers \(I\), on peut poser \(x=\varphi (t)\), \(x\in I\), on a :
	            \[
	                  \int f(x) ~\mathrm{d}x = \int f(\varphi(t)) \cdot \varphi'(t) \mathrm{d}t      \,.
	            \]
\end{frame}
\begin{frame}[allowframebreaks]{Intégrales (2)}
example
	                  Posons \(f(x) = x \sqrt{x+1}\).
	                  On veut trouver une primitive de \(f\).
	                  Posons \(t = \sqrt{x+1}\), on a \(x = t^2 - 1\) et donc \(~\mathrm{d}x = 2t \mathrm{d}t\).
	                  \[\begin{aligned}
	                        \int x \sqrt{x+1} ~\mathrm{d}x & = \int (t^2 - 1) \cdot t \cdot 2t \mathrm{d}t                                     
	                                                = \int (2t^4 - 2t^2) \mathrm{d}t                                                  \\
	                                              & = 2 \left[ \dfrac{t^5}{5} - \dfrac{t^3}{3} \right] + C                      
	                                                = \dfrac{2}{5} (x+1)^{\dfrac{5}{2}} - \dfrac{2}{3} (x+1)^{\dfrac{3}{2}} + C   \,.
	                  \end{aligned}\]
\end{frame}
\begin{frame}[allowframebreaks]{Intégrales (3)}
Ainsi, la fonction \(F\) définie par 
	                  \(F(x) = \dfrac{2}{5} (x+1)^{\dfrac{5}{2}} - \dfrac{2}{3} (x+1)^{\dfrac{3}{2}}\)
	                  est une primitive de \(f\) sur \(]-1,+\infty[\) (on revient à la variable \(x\) grâce à \(t = \sqrt{x+1}\)).
	            example
\end{frame}
\begin{frame}[allowframebreaks]{Intégrales}
\footnotesize
Soient \(f\) une fonction continue et positive sur un intervalle \([a,b]\) (\(a<b\))
      et notons par \(\mathcal{C}\) sa courbe représentative dans 
      le plan rapporté à un repère orthogonal \((O,\overrightarrow{\imath},\overrightarrow{\jmath})\).
      L'aire du domaine limité par \(\mathcal{C}\), l'axe \(x'Ox\) et les droites d'équations \(x=a\) et \(x=b\)
      est alors égal à \(\int_{a}^{b} f(x) ~\mathrm{d}x\).
      Cette valeur s'exprime en unités d'aire : l'unité d'aire est l'aire du rectangle construit sur \(\overrightarrow{\imath}\) et \(\overrightarrow{\jmath}\) (de côtés \(\|\overrightarrow{\imath}\|\) et \(\|\overrightarrow{\jmath}\|\)).
\end{frame}
\begin{frame}[allowframebreaks]{Intégrales}
Notons que si \(f\) est négative sur \([a,b]\), alors la valeur de \(\int_{a}^{b} f(x) ~\mathrm{d}x\) est aussi négative.

Ainsi, dans ce cas, l'aire délimitée par \(\mathcal{C}\), l'axe \(x'Ox\) et les droites d'équations \(x=a\) et \(x=b\)
      est alors égal à \(-\int_{a}^{b} f(x) ~\mathrm{d}x\).
\end{frame}
\begin{frame}[allowframebreaks]{Intégrales}
\footnotesize
Plus généralement, si la fonction prend à la fois des valeurs positives et négatives sur \([a,b]\),
      l'aire du domaine plan correspond à \(\int_{a}^{b} |f(x)| ~\mathrm{d}x\).
      La valeur absolue permet de s'assurer qu'on calcule toujours des valeurs positives.
      Pour calculer l'aire du domaine limité par la courbe d'une fonction qui change de signe,
      on peut diviser l'intervalle considéré en plusieurs sous-intervalles
      de manière à ce que le signe de la fonction soit constant sur chaque sous-intervalle.
      On considère alors la valeur absolue des surfaces ainsi calculées.
      Il suffit ensuite d'additionner les aires trouvées pour obtenir l'aire totale.
      Si on calcule juste la valeur de \(\int_{a}^{b} f(x) ~\mathrm{d}x\), alors 
      la partie négative et la partie positive se soustrairont.
\end{frame}
\begin{frame}[allowframebreaks]{Intégrales}
Soient \(a>0\) et \(f\) une fonction continue sur \([-a,a]\).

\begin{itemize}

                  \item Si \(f\) est une fonction paire alors \(\int_{-a}^{a} f(x) ~\mathrm{d}x = 2 \cdot \int_{0}^{a} f(x) ~\mathrm{d}x\).
                  
\item Si \(f\) est une fonction impaire alors \(\int_{-a}^{a} f(x) ~\mathrm{d}x = 0\).

\end{itemize}
\end{frame}
\begin{frame}[allowframebreaks]{Intégrales}
Soit \(f\) une fonction continue sur \([a,b]\) et \(\psi : J \rightarrow [a,b]\) une fonction dérivable, de dérivée continue, et bijective.
      
      Si \(\psi (\alpha) = a\) et \(\psi (\beta) = b\), alors : 
            \[
                  \int_{a}^{b} f(x) ~\mathrm{d}x = \int_{\alpha}^{\beta} f(\psi(t)) \cdot \psi'(t) \mathrm{d}t      \,.
            \]
\end{frame}
\begin{frame}[allowframebreaks]{Intégrales (2)}
example
                  Posons \(f(x) = x \sqrt{x+1}\).
                  On veut trouver l'integrale \(\int_{0}^{1} f(x) ~\mathrm{d}x\).
                  Posons \(t = \sqrt{x+1}\), on a \(x = t^2 - 1\) et donc \(~\mathrm{d}x = 2t \mathrm{d}t\).
                  Si \(x=0\) alors \(t=1\) et si \(x=1\) alors \(t=\sqrt{2}\).
                  \[\begin{aligned}
                  \int_{0}^{1} x \sqrt{x+1} ~\mathrm{d}x & = \int_{1}^{\sqrt{2}} (t^2 - 1) \cdot t \cdot 2t \mathrm{d}t                      
                                                  = \int_{1}^{\sqrt{2}} (2t^4 - 2t^2) \mathrm{d}t                                   \\
                                                & = 2 \left[ \dfrac{t^5}{5} - \dfrac{t^3}{3} \right]_{1}^{\sqrt{2}}           
                                                  = 2 \left[ 
                                                      \left( \dfrac{4\sqrt{2}}{5}-\dfrac{2\sqrt{2}}{3} \right) 
                                                      -
                                                      \left( \dfrac{1}{5} - \dfrac{1}{3} \right)
                                                      \right]                                                               \\
                                                & = \dfrac{4}{15} \cdot \left(\sqrt{2} + 1 \right)                           \,.
                  \end{aligned}\]
\end{frame}
\begin{frame}[allowframebreaks]{Intégrales (3)}
example
\end{frame}
\begin{frame}[allowframebreaks]{Intégrales}
Soit \(f\) une fonction définie et continue sur un intervalle \(I\).
\end{frame}
\begin{frame}[allowframebreaks]{Intégrales}
De la forme \(\int_{a}^{x} f(t) \mathrm{d}t\)

 Comme \(f\) est continue, elle admet des primitives sur \(I\).
                  Pour tout \(a \in I\), on peut définir la fonction \(F\) par \(F(x) = \int_{a}^{x} f(t) \mathrm{d}t\),
                  qui n'est autre que la primitive de \(f\) qui s'annule en \(a\).
                  Cette fonction est dérivable et a pour dérivée \(f\).
\end{frame}
\begin{frame}[allowframebreaks]{Intégrales}
De la forme \(\int_{u(x)}^{v(x)} f(t) \mathrm{d}t\)

 Soient \(u\) et \(v\) deux fonctions dérivables dans un domaine \(\mathcal{D}\), et à valeurs dans \(I\).
                  La fonction \(f\) étant continue, elle admet une primitive \(F\)
                  En fait elle en admet plusieurs, on va juste considérer une d'elles.
\end{frame}
\begin{frame}[allowframebreaks]{Intégrales (2)}
On peut alors définir la fonction \(H\) par 
                  \(H(x) = \int_{u(x)}^{v(x)} f(t) \mathrm{d}t\).
                  On sait que 
                  \[\begin{aligned}
                        H(x) & = \int_{u(x)}^{v(x)} f(t) \mathrm{d}t = F(v(x)) - F(u(x)) \;.
                  \end{aligned}\]

Cette fonction est dérivable sur \(\mathcal{D}\) et sa dérivée est définie par

                              \[
                                    H'(x) = v'(x) \cdot f(v(x)) - u'(x) \cdot f(u(x)) .
                              \]
\end{frame}
\end{document}
